\documentclass[11pt]{article}
\usepackage[a4paper,margin=1in]{geometry}
\usepackage{fontspec}
\usepackage[boldfont=false]{xeCJK}
\setCJKmainfont{FandolSong-Regular.otf}
\usepackage{amsmath}
\usepackage{amssymb}
\usepackage{array}
\usepackage{booktabs}
\usepackage{graphicx}
\usepackage{adjustbox}
\usepackage{enumitem}
\setlist[itemize]{topsep=2pt,partopsep=0pt,parsep=0pt,itemsep=2pt,leftmargin=1.4em}
\setlist[enumerate]{topsep=2pt,partopsep=0pt,parsep=0pt,itemsep=2pt,leftmargin=1.6em}
\usepackage{parskip}
\usepackage{fvextra}
\fvset{breaklines=true,breakanywhere=true,fontsize=\small}
\setlength{\parindent}{0pt}

\begin{document}

\begin{center}
  {\LARGE\bfseries Hydrocarbons}\\[0.35em]
  {\large A Level Chemistry}
\end{center}
\vspace{0.6em}

\section*{Arenes}

\textbf{Arenes} 芳烃 are aromatic hydrocarbons, built on the \textbf{benzene} 苯 ring. The delocalised ring of electrons is stable and electron-rich, so benzene mostly reacts by \textbf{electrophilic substitution} 亲电取代 — keeping the ring — rather than by addition.

\subsection*{Reactions of benzene and methylbenzene}

\begin{center}
\adjustbox{max width=\linewidth}{%
\begin{tabular}{lll}
\toprule
\textbf{Reaction} & \textbf{Reagents and conditions} & \textbf{Product} \\
\midrule
halogenation & $\text{Cl}_2$ or $\text{Br}_2$, with $\text{AlCl}_3$ or $\text{AlBr}_3$ as a \textbf{catalyst} 催化剂 & a \textbf{halogenoarene} 卤代芳烃 (aryl halide) \\
\textbf{nitration} 硝化 & concentrated $\text{HNO}_3$ + concentrated $\text{H}_2\text{SO}_4$, $25$–$60\,°\text{C}$ & nitrobenzene \\
\textbf{Friedel–Crafts alkylation} 傅克烷基化 & $\text{CH}_3\text{Cl}$ + $\text{AlCl}_3$, heat & methylbenzene (adds an alkyl group) \\
\textbf{Friedel–Crafts acylation} 傅克酰基化 & $\text{CH}_3\text{COCl}$ + $\text{AlCl}_3$, heat & a phenyl ketone (adds an acyl group) \\
side-chain oxidation & hot alkaline $\text{KMnO}_4$, then dilute acid & \textbf{benzoic acid} 苯甲酸 \\
\textbf{hydrogenation} 氢化 & $\text{H}_2$, $\text{Pt/Ni}$ catalyst, heat & cyclohexane \\
\bottomrule
\end{tabular}}
\end{center}

In the side-chain oxidation, the whole \textbf{side-chain} 侧链 (such as the $\text{–CH}_3$ on methylbenzene) is turned into a $\text{–COOH}$ group, giving benzoic acid. In the hydrogenation, three molecules of $\text{H}_2$ add to the \textbf{benzene ring} 苯环 to make a saturated cyclohexane ring.

\subsection*{The mechanism: electrophilic substitution}

Take nitration as the example. The acid mix makes the electrophile $\text{NO}_2^+$. Then:

\begin{enumerate}
\item the delocalised electrons of the ring form a bond to the electrophile, giving an unstable intermediate.

\item an $\text{H}^+$ is lost from that carbon, which restores the stable ring.

\end{enumerate}

Substitution wins over addition because the \textbf{delocalisation} 离域 (aromatic stabilisation) of the ring is kept. Addition would destroy this stable system, so it is not favoured.

\subsection*{Side-chain or ring?}

Where a halogen reacts depends on the conditions:

\begin{itemize}
\item with a halogen-carrier catalyst (such as $\text{AlCl}_3$) and \textbf{no} UV light → substitution in the \textbf{ring}.

\item with UV light and \textbf{no} catalyst → free-radical substitution in the \textbf{side-chain}.

\end{itemize}

\subsection*{Directing effects}

A group already on the ring decides where the next group goes — its \textbf{directing effect} 定位效应:

\begin{center}
\adjustbox{max width=\linewidth}{%
\begin{tabular}{ll}
\toprule
\textbf{Group already present} & \textbf{Directs the new group to} \\
\midrule
$\text{–NH}_2$, $\text{–OH}$, $\text{–R}$ & positions 2 and 4 \\
$\text{–NO}_2$, $\text{–COOH}$, $\text{–COR}$ & position 3 \\
\bottomrule
\end{tabular}}
\end{center}


\end{document}
